\section{The Muzzle of a Pistol}

A review of \emph{Against Nature} by \emph{Joris-Karl Huysmans}, in which the decadence of the pursuit of phenomena is exposed.

\begin{quotex}
After such a book, the only choice left open to the author is the muzzle of a pistol or the foot of the cross. \flright{From Barbey d'Aurevilly's review of \emph{Against Nature} by \textbf{Joris-Karl Huysmans}}

Seeing that nowadays there is nothing wholesome left in this world of ours; seeing that the wine we drink and the freedom we enjoy are equally adulterate and derisory; and finally, seeing that it takes a considerable degree of goodwill to believe that the governing classes are worthy of respect and that the lower classes are worthy of help or pity, it seems to me no more absurd or insane to ask of my fellow men a sum total of illusion barely equivalent to that which they expend every day on idiotic objects. \flright{from \emph{Against Nature}}

Verily it is a pitiable thing to live on earth \flright{from the \emph{Imitation of Christ}}

\end{quotex}
\textbf{Joris-Karl Huysmans} was present at the end of Naturalism in art. By attempting to place so-called ``real characters" in ``precise settings", the Naturalists ended up just repeating themselves. In a preface published 20 years after the initial publication of \emph{Against Nature} (\emph{A Rebours}) in 1884, Huysmans explained that since Virtue was such an exception on earth, it was excluded from the Naturalist framework. He wrote:

\begin{quotex}
It must be admitted that Virtue is an exception here on earth, and so it was excluded from the Naturalist framework. Not having the Catholic conception of the fall from grace and of temptation, we were unaware of the struggles and sufferings from which Virtue springs; the heroism of the soul, triumphing over life's pitfalls, escaped us. It would not have occurred to us to describe this struggle, with its highs and lows, its wily attacks and feints, or the able allies standing at the ready deep in their cloisters and often far from those the Devil is assailing. Virtue seemed to us the attribute of individuals without curiosity or bereft of sense, and in any case of too little emotional interest to treat from an artistic point of view. 

\end{quotex}
Without the spiritual combat of virtue against vice, the Naturalists could only write about the latter. Of course, the only interesting vice turns out to be Lust. Eventually, even that topic gets played out, at least for the intelligent artistic mind. Nevertheless, in popular culture it remains a popular theme, since the mass mind cannot even conceive of an alternative.

\paragraph{Addiction to Phenomena}
If naturalism is true, then only what can be experienced by the senses exists: sight, sound, touch, taste, scent. Hence, the goal of life is to accumulate the best experiences possible. For most, this is mostly a passive way of life: attending concerts, choosing a tattoo, finding a great restaurant, watching sports, and the like. A more sophisticated person will have refined tastes. For example, a woman recently wrote to me with her goals for the new year:

\begin{quotex}
Freedom, authenticity, creativity, love, connection… laughter, fat books, Eddie Higgins, eating with the locals…. oysters, avocados, ripasso, chocolate with sea salt, rose-scented soap. 

\end{quotex}
A cute list, but still rather passive. The novel \emph{Against Nature} is a story about nothing, except the pursuit of sensual experiences. The sole character, Des Esseintes, is a recluse with the resources in money and time to indulge his tastes. The difference is that he is the creator of his experiences. In this review, we can only highlight some of these.

First, he is obsessed with the way the changing light in his room reflects the multicolored threads of his rug. Whatever he does involves intricate attention to detail. To better reflect the sunlight in his room, he goldplates a tortoise and covers the shell with carefully selected gemstones. Needless to say, the tortoise died, probably due to a disease of the shell. He built a ``mouth organ", which controlled a variety of various liqueurs of various colors and shapes. That invention was reprised as the pianocktail in the French film comedy \textit{Mood Indigo}\footnote{\url{https://www.youtube.com/watch?v=5bdX5i0nAWw}}.

He had an extensive library, with many titles privately and extravagantly printed. They were arranged like the Louvre, from the ancients to the medievals, and ending with contemporary French literature. Curiously, it included an extensive collection of religious texts although Des Esseintes was an unbeliever. Some of it was of good quality and some merely sentimental. Apparently, he had been quite experimental with sex in the Parisian whorehouses. This left him impotent, probably because conventional sex could no longer arouse him.

Despite his further experiences with exotic flowers, literature, and fine art, he remained dissatisfied. Eventually, his health began to deteriorate. His physician recommended that Des Esseintes return to live in the city and ``try to enjoy the same pleasures as other people."

The thought of that was repulsive. He prided himself on his ability to appreciate ``the delicacy of a phrase, the subtlety of a painting, the quintessence of an idea, or whose soul was sensitive enough to understand Mallarme and love Varlaine." Where could he find a ``twin soul":

\begin{quotex}
a mind free of commonplace ideas, welcoming silence as a boon, ingratitude as a relief, suspicion as a haven and a harbour. 

\end{quotex}
There it is: the pursuit of experiences for their own sake leads to death, not only of the body, but especially of the soul. The only possible solution is to transcend all experience, yet he couldn't bring himself to the Faith. Hence, he expects to be engulfed by ``waves of human mediocrity". With no hope of a future life, his earthly life is all there is. Nevertheless, the atheist ends with a prayer:

\begin{quotex}
Lord, take pity on the Christian who doubts, on the unbeliever who would fain believe, on the galley-slave of life who puts out to sea alone, in the night, beneath a firmament no longer lit by the consoling beacon-fires of the ancient hope. 

\end{quotex}

\paragraph{The Foot of the Cross}

Here we review Huysmans' preface for the edition of \emph{A Rebours} published 20 years later. In particular, it deals with the influence of \textbf{Arthur Schopenhauer} and \textbf{John of Ruysbroeck} which led Huysmans to take the vows of a Cistercian oblate.


\hfill

Joris-Karl Huysmans describes his journey from Des Esseints to the Cistercians. 

\begin{quotex}
I must rejoice beyond the bounds of time … though the world may shudder at my joy, and in its coarseness know not what I mean \flright{\textsc{John of Ruysbroeck}}

John of Ruysbroeck was a thirteenth century mystic whose prose presented an incomprehensible but attractive amalgam of gloomy ecstasies, tender raptures and violent rages. \flright{Des Esseintes, in \emph{Against Nature}}

Extraordinary things can only be stammered out, and stammer [John of Ruysbroeck] did, declaring that the sacred obscurity in which Ruysbroeck spreads his eagle's wings is his ocean, his prey, his glory, and for him the four horizons would be too close-fitting a garment. \flright{\textsc{Ernest Hello}}

There is more knowledge and understanding of the human heart in one page of Ruysbroeck than in all the Stendhals, Bourgets, and Barreses in the world. \flright{\textsc{Joris-Karl Huysmans}}

\end{quotex}
In a new Preface to the 20th anniversary edition of \emph{Against Nature}, \textbf{Joris-Karl Huysmans} tried to explain how, eight years after its original publication, he ended up taking vows in a Cistercian monastery. However, the full story of his conversion was related in three subsequent books, we will focus at this time solely on what he wrote in the Preface. In particular, what is of interest is the influence of two writers, one a philosopher, the other a mystic, who achieved a certain prominence at that time. The philosopher was \textbf{Arthur Schopenhauer}, who was achieving a posthumous reputation that had eluded him in his lifetime. The mystic was \textbf{John of Ruysbroek} whose \emph{Adornment of the Spiritual Marriage} had recently been translated into French by the theologian Ernest Hello.

After posting some extracts from Husymans' Preface, we will conclude later this week with some fascinating ideas from Ruysbroeck. These will show how differently doctrine is understood today compared to the Medieval era. The important point to notice is to observe how a great artist approaches things. Like the poet, Huysmans seemed to have had negative capability and was willing to consider many ideas, while avoiding simple answers. Nowadays, in the age of social media, thoughts that require more than a dozen words are ignored. Entire worldviews are expressed in seven declarative sentences, often with a florid background. Negative views are viewed negatively.

The idea of living with uncertainty and mystery is terrifying; all must be known and under control. Huysmans wrote:

\begin{quotex}
I did not grasp that all is mystery, that we live only in mystery, that if such a thing as chance existed it would be even more mysterious than Providence. 

\end{quotex}
The lack of negative capability makes an artistic conversion like Huysmans' rare today. For example, what unbeliever could make this observation:

\begin{quotex}
The strange thing was that, without my realizing it at first, I was drawn by the nature of my work itself to study the Church from a number of angles. It was in fact impossible to trace one's way back to the only unblemished times humanity has ever known, the Middle Ages, without realizing that the Church was at the centre of everything, that art existed through her and by her. Not being a believer, I looked at her, a little defiant, taken aback by her greatness and her glory, wondering why a religion which seemed to me to have been created for children could have inspired such marvelous works of art. 

\end{quotex}
At that time, Huysmans was under Schopenhauer's spell, who was the philosopher of the tragic side of life. Like Huysmans, Schopenhauer was aware of religious literature like the Upanishads, Buddha, and even Christian mysticism. They all acknowledged what Schopenhauer observed, except that they offered a solution, an escape.

\begin{quotex}
Yet I thought myself so far from religion. I did not imagine that it was only a short step from Schopenhauer, whom I admired beyond reason, to \emph{Ecclesiastes} and the \emph{Book of Job}. The premises about Pessimism are the same, only when the time comes to reach a conclusion, the philosopher disappears. I like his ideas about the horror of life, the stupidity of the world, the mercilessness of destiny; I like them also in the Holy Scriptures; but Schopenhauer's observations lead nowhere; he leaves you, so to speak, in the lurch; in the end, his aphorisms are only a herbarium of dry plaints; whereas the Church explains the origins and the causes, indicates the conclusions, offers remedies. She does not limit herself to giving you a spiritual consultation, but treats and cures you, whereas the German quack, once he has proved the incurability of your condition, simply sneers and turns his back on you. … Never has Pessimism been of any comfort to those sick in body or in soul! 

\end{quotex}
There was no single event that led Huysmans to his conversion. He tries to explain the inexplicable:

\begin{quotex}
There was undoubtedly, as I was writing \emph{Against Nature}, a land-shift, the earth was being mined to lay foundations of which I was unaware. God was digging to set his fuses and he worked only in the darkness of the soul, in the night. Nothing could be seen; it was only years later that the sparks began to run along the wires. I felt my soul moving to these shocks; it was at the time neither especially painful nor especially clear. … During the liturgies, I felt nothing more than an inner trepidation, a trembling that one feels when one sees or hears or reads a beautiful work of art, but there was no precise warning to get me ready to make up my mind. 

\end{quotex}
He continues:

\begin{quotex}
I was simply emerging, little by little, from the shell of my moral impurity; I was beginning to be disgusted with myself, but still I balked at the articles of faith. The objections I placed in my path seemed irresistible to me; and one morning, when I awoke, they were resolved, how, I have never known. I prayed for the first time and the explosion happened. 

\end{quotex}
In these events, Huysmans was following the course described by \textbf{John of Ruysbroeck}:

\begin{quotex}
And from this love there come forth perfect contrition and purification of conscience. And these arise from the consideration of misdeeds and all that may defile the soul: for when a man loves God he despises himself and all his works. 

\end{quotex}
Selected texts from John of Ruysbroeck: \url{https://gornahoor.net/?p=8452}


\hfill



\flrightit{Posted on 2016-01-06 by Cologero }

\begin{center}* * *\end{center}

\begin{footnotesize}\begin{sffamily}



\texttt{winchelljames1 on 2020-01-08 at 23:40 said: }

Your opening lines seem to imply so much by saying so little, yet the little information they DO manage to convey ignores so much substance and nuance as to qualify as a kind of vacuous ``negative space" of both literary criticism and history, to wit: ``Joris-Karl Huysmans was present at the end of Naturalism in art." Well, yes, this is correct as far as it goes–but as a young disciple of Emile Zola, who attended the regular salons of the Naturalist Master at his home in Médan, Huysmans was even more remarkably present–and active–at the coming-of-age of Naturalism as a literary-artistic movement, when those disciples' early prose works, written under Zola's tutelage and edited collectively as *Les Soirées de Médan,* were published in 1880–the same year Zola published his Naturalist manifesto, ``Le Roman expérimental." Such demurrers may be dismissed as ``the tall grass" by those not in the field, but the truly off-handed overstatement and dismissive attitude conveyed in your second sentence may raise bristling hackles on the necks of even the most determined generalist: ``By attempting to place so-called `real characters' in `precise settings', the Naturalists ended up just repeating themselves." This second clause does a gross disservice-by-reductionism to Zola's theory and practice of Naturalism as he developed it over a span of some 20 novels (and they were big, ambitious, bulky, door-stop volumes, mind you) gathered under the names of the two families whose genealogies he attempted to trace (*Les Rougon-Macquart*), AND as he theorized the movement in the above-mentioned essay (``Le Roman expérimental"), published in 1880. The latter proves to be, upon inspection by anyone who might actually bother to read it, an extremely rational, detailed, pioneering (if only quasi-scientific, if not merely–and enjoyably–``Quack-xotic") synthesis of then-current theories of evolution as had come down to him from Hippolyte Taine and Auguste Comte. Your gloss on the authors' commitment to Realism at the heart of Naturalism sounds both musty and snobby (``so-called `real characters'…), but I take most offense at your final clause: ``…the Naturalists ended up just repeating themselves." This complacent, dismissive, cocktail-party apothegm may sound informed and succinct, but alas reveals a most basic failure to grasp even the most basic premise of the movement: by recreating in the petri-dishes of fiction the evolutionary life-forces of aberrant breeding (``race" en français), specific history (``moment") and recognizable locale (``milieu," topos), Zola (naively, of course) believed that he might trace–or even predict–one, or some, of the evolutionary traits in the biology of human beings. You may be mistaking the ``repetitions" you so blithely dismiss for the background of non-aberrations occurring ``naturally" within any specimen-sample, or, of course, in ``Nature" herself. Of course–Modern Art, or the advent of Impressionism and all subsequent Abstractions, begins with nothing less than the statement: Seen one landscape, seen 'em all! Is it merely an accident–coincidence?–that Impressionism is flowering at pretty much the same time as Naturalism is running–and completing–its course?


\hfill

\texttt{Cologero on 2020-01-09 at 08:18 said: }

By the way, what you attributed to me was actually a quote from Huysmans' preface to the novel. So you are hardly the expert you seem to believe you are. From Husmans: ``In those days Naturalism was at its height; but this school of writers, which was to fulfil the invaluable service of placing real characters in precise settings, was condemned to repeat itself over and over, and endlessly to go over the same ground."

Which Zola himself did in his 20 prolix volumes, endlessly repeating himself.


\hfill

\texttt{winchelljames1 on 2020-01-10 at 02:48 said: }

Au contraire, mon frère.

Cf. my complete response on FB page ``Christian Pessimist Thinkers."


\hfill

\texttt{santiago on 2020-01-10 at 21:17 said: }

I know next to nothing about the late history of Naturalism in art, so I hope this post won't be entirely irrelevant, but; in my experience, the `life is a theme park' mentality is especially dangerous for those who have a `transcendent', `adventurous' or `pioneering' inner constitution. Among these types, when infected with this disease of `experience collecting' or `thrill-seeking', one can find drug addicts, gamblers, sexual degenerates, brawlers, adrenaline junkies, etc. Re-reading Nietzsche recently, I think he, better than most, exemplified this tendency of a drive towards experience and self-actualization, and the consequences thereof, when divorced from any higher principles. It really is a shame. The darkest moments of my life have usually been preceded by having adopted this same mentality beforehand. I learned my lesson.

This reduction of life to a collection of separate `experiences' only significant in themselves, and only measurable in their `authenticity' as decided by the personal subjective observer, has consequences far greater than just addiction, depression and ambivalence. ``By absurd and destructive actions, they sought to attain the only possible meaning of existence, after rejecting suicide as the radical solution for the metaphysically abandoned individual." 

``Needless to say, the tortoise died, probably due to a disease of the shell."


\end{sffamily}\end{footnotesize}
